% Topic presentation, 20 minutes. Template Beamer dimensions and theme retained.
\documentclass[aspectratio=169,11pt]{beamer}
\usepackage[utf8]{inputenc}
\usepackage[T1]{fontenc}
\usepackage[spanish,es-nodecimaldot]{babel}
\usepackage{amsmath,amssymb,booktabs}
\usepackage{pgfplots}
\usepackage{listings}
\pgfplotsset{compat=1.18}
\usetheme{default}
\usecolortheme{seahorse}
\setbeamertemplate{navigation symbols}{}
\setbeamertemplate{footline}[frame number]
\setbeamerfont{frametitle}{size=\large}
\newcommand{\E}{\mathbb E}
\newcommand{\Var}{\operatorname{Var}}
\lstset{basicstyle=\ttfamily\scriptsize,columns=fullflexible,
  keepspaces=true,breaklines=false,showstringspaces=false,
  escapeinside={(*@}{@*)},
  literate={ℝ}{{\(\mathbb R\)}}1 {≤}{{\(\leq\)}}1}
\title{¿Cuándo la transparencia mejora la\\incorporación de fundamentos al precio?}
\subtitle{Presentación del tema, ruta A}
\author{David Julio Olano Silva Nevado}
\date{9 de octubre de 2026}

\begin{document}
\begin{frame}[plain]
 \titlepage
 \begin{center}\small
 \url{https://github.com/davidolano03/ai-project}
 \end{center}
\end{frame}

\begin{frame}{1. Pregunta de investigación}
 \large
 ¿Bajo qué condiciones una mayor transparencia aumenta la
 \textbf{sensibilidad esperada del precio a los fundamentos},
 cuando algunos inversionistas están afectados por sentimiento?
 \vspace{0.6cm}
 \normalsize
 \begin{itemize}
  \item Mediré cuánto responde el precio a una misma noticia fundamental.
  \item Una respuesta más informativa puede elevar o reducir el precio.
  \item Corregiré la información parcial de Zhang y Zhang (2023).
 \end{itemize}
\end{frame}

\begin{frame}{2. El supuesto que cambiaré}
 Zhang y Zhang (2023), sección 6.2, ecuaciones (38)--(41), pp.~15--16:
 \[
   \kappa=\mu\theta,\qquad
   v_3^I=\sigma_\epsilon^2+\mu^2\sigma_\theta^2.
 \]
 \begin{itemize}
  \item $\theta$: noticia fundamental. $\kappa$: información pública.
        $\mu\in[0,1]$: transparencia.
  \item Con $\mu>0$ conocido, observar $\kappa$ revela $\theta=\kappa/\mu$.
  \item En $\mu=0$, la varianza impresa elimina la incertidumbre
        sobre una noticia que todavía no se observa.
 \end{itemize}
 \medskip
 \textbf{Resultado afectado:} proposición 5(ii), p.~16.
 Afirma que todo grado de transparencia mejora la sensibilidad informativa.
\end{frame}

\begin{frame}{Evidencia preliminar del problema}
 Las fórmulas publicadas (23) y (44) admiten este caso:
 \begin{center}
 \begin{tabular}{lc}
 \toprule
 Escenario & Sensibilidad informativa\\
 \midrule
 Modelo básico & $9/14=0{,}642857$\\
 Transparencia publicada, $\mu=0{,}1$ & $66943/121303=0{,}551866$\\
 \bottomrule
 \end{tabular}
 \end{center}
 \[
  0{,}551866<0{,}642857.
 \]
 \small
 $\lambda_i=1/4$, $R=\gamma=1$, sentimientos y medias cero.
 Todas las varianzas del modelo básico son uno.
 \medskip
 \normalsize
 Lean verifica el contraejemplo. La afirmación universal falla
 \textbf{con las fórmulas impresas}; esto no establece un efecto empírico.
\end{frame}

\begin{frame}{Información pública con incertidumbre residual}
 Propondré componentes gaussianos independientes, de media cero:
 \[
   \theta=\kappa+\xi,\qquad
   \Var(\kappa)=\mu s,\qquad \Var(\xi)=(1-\mu)s,
   \quad s=\sigma_\theta^2.
 \]
 Al observar $\kappa$, el particular tiene
 \[
   \E[\theta\mid\kappa]=\kappa,\qquad
   \Var(\theta\mid\kappa)=(1-\mu)s.
 \]
 \begin{itemize}
  \item $\mu$ es una fracción de varianza divulgada.
  \item Más información reduce la incertidumbre restante.
  \item $\mu=0$ recupera el modelo básico.
        $\mu=1$ elimina el riesgo informativo, pero no la sorpresa $\epsilon$.
 \end{itemize}
\end{frame}

\begin{frame}{Los cuatro grupos del modelo corregido}
 \small
 $F=P_0+\theta+\epsilon$, con $P_0$ precio previo;
 $\epsilon$ independiente, de media cero y varianza $\sigma_\epsilon^2>0$.
 \begin{center}
 \begin{tabular}{@{}llll@{}}
 \toprule
 Grupo & Media $m_i$ & Varianza $v_i$ & Aversión\\
 \midrule
 Institución racional & $P_0+\theta$ & $\sigma_\epsilon^2$ & $\gamma$\\
 Institución con sentimiento & $P_0+\theta+\delta_1$ & $\sigma_\epsilon^2+\tau_1^2$ & $g_2$\\
 Particular racional & $P_0+\kappa$ & $\sigma_\epsilon^2+(1-\mu)s$ & $\gamma$\\
 Particular con sentimiento & $P_0+\kappa+\delta_2$ & $\sigma_\epsilon^2+(1-\mu)s+\tau_2^2$ & $g_4$\\
 \bottomrule
 \end{tabular}
 \end{center}
 \[
  g_2=\gamma e^{-b\rho_1},\qquad g_4=\gamma e^{-b\rho_2},
  \qquad b,\gamma>0.
 \]
 \normalsize
 $\delta_j=\rho_j^*$: sesgo esperado.
 $\tau_j^2=\sigma_{\rho_j}^2$: incertidumbre subjetiva adicional.
 $\rho_j$: estado emocional observado.
 \medskip
 \par
 Conservaré los cuatro grupos. Cambiaré las creencias informativas
 de los particulares.
\end{frame}

\begin{frame}{3. Decisión y FOC esperada}
 El agente elige su posición $\omega_i\in\mathbb R$:
 \[
  \max_{\omega_i}\;
  C_{0,i}+\omega_i(m_i-RP)-\frac{g_iv_i}{2}\omega_i^2,
  \qquad R=1+r_f>0.
 \]
 \[
   m_i-RP-g_iv_i\omega_i=0,\qquad
   \omega_i^*=\frac{m_i-RP}{g_iv_i}.
 \]
 \begin{itemize}
  \item $C_{0,i}$: riqueza independiente de la decisión.
  \item Ganancia esperada menos una penalización por riesgo.
  \item Para el particular racional, $v_3=\sigma_\epsilon^2+(1-\mu)s$.
 \end{itemize}
 Con precio y expectativas iguales, menor incertidumbre amplía la posición:
 más compra si es positiva y más venta corta si es negativa.
\end{frame}

\begin{frame}{Precio y relevancia de cada grupo}
 Con participaciones $\lambda_i>0$, $\sum_i\lambda_i=1$ y oferta de una unidad:
 \[
  \sum_i\lambda_i\omega_i^*=1,\qquad
  h_i=\frac{\lambda_i}{g_iv_i},\qquad H=\sum_i h_i,
 \]
 \[
   P^*=\frac{\sum_i h_im_i-1}{RH}.
 \]
 \begin{itemize}
  \item $h_i$ mide el peso efectivo de demanda.
  \item Más participación, menos aversión o menos riesgo elevan ese peso.
  \item Si $m_i>RP$, darle más peso al grupo eleva el precio.
        Si $m_i<RP$, lo reduce.
 \end{itemize}
 La relevancia de un grupo tiene un efecto condicionado por su valoración.
\end{frame}

\begin{frame}{Resultado esperado y condiciones}
 Mediré la respuesta del \textbf{precio esperado} a un mismo fundamento:
 \[
  \bar P(\theta,\mu)=\E[P^*\mid\theta,\rho_1,\rho_2],
  \qquad \E[\kappa\mid\theta]=\mu\theta.
 \]
 \[
  \mathcal K(\mu)=\frac{\partial\bar P}{\partial\theta}
    =\frac{a+\mu B(\mu)}{R[a+B(\mu)]},
  \quad a=h_1+h_2,\quad B=h_3+h_4.
 \]
 \[
  0\leq x<y\leq1
  \quad\Longrightarrow\quad
  \mathcal K(y)>\mathcal K(x),
  \qquad \mathcal K(1)=1/R.
 \]
 \small
 Participaciones positivas, $R>0$, $\sigma_\epsilon^2>0$,
 $s>0$, $\tau_j^2\geq0$. Sentimiento exógeno e independiente de $\kappa,\xi$.
 Mantendré fijos sus estados, sesgos, participaciones y varianzas totales.
 \medskip
 \normalsize
 Más transparencia debería fortalecer la respuesta a buenas
 \textbf{y} malas noticias.
\end{frame}

\begin{frame}{Comparación reproducible de la sensibilidad}
 \begin{center}
 % Generated from checked data. No screenshots of published figures.
\begin{tikzpicture}
\begin{axis}[
 width=10.6cm,height=5.7cm,
 xmin=0,xmax=1,ymin=0.45,ymax=1.03,
 xlabel={Transparencia $\mu$},ylabel={Sensibilidad informativa $\mathcal K$},
 tick label style={font=\small},label style={font=\small},
 legend style={font=\scriptsize,draw=none,at={(0.02,0.98)},anchor=north west},
 grid=major,grid style={gray!15},
]
\addplot[blue!70!black,very thick] table[x=mu,y=corrected,col sep=comma]{../code/output/information.csv};
\addlegendentry{Corrección: precio esperado}
\addplot[orange!85!black,thick] table[x=mu,y=printed,col sep=comma]{../code/output/information.csv};
\addlegendentry{Fórmula publicada}
\addplot[black,dashed] table[x=mu,y=baseline,col sep=comma]{../code/output/information.csv};
\addlegendentry{Modelo básico}
\end{axis}
\end{tikzpicture}

 \end{center}
 \vspace{-0.3cm}
 \scriptsize
 Parámetros del contraejemplo. La curva corregida mide respuesta esperada.
 Datos generados por \texttt{code/verify.py}, no evidencia empírica.
\end{frame}

\begin{frame}{4. Por qué esta corrección requiere trabajo propio}
 \begin{itemize}
  \item \textbf{Paper y apéndice:} Zhang ya incorpora aprendizaje y
        transparencia. La proposición 5 mantiene la parametrización discutida.
  \item \textbf{Versiones de trabajo:} búsqueda por título y DOI.
        No identifiqué una versión alternativa accesible ni una corrección.
  \item \textbf{Artículo citante revisado:} Das y Yaghoubi (2024) estudia
        miedo y riesgo bursátil empíricamente. Usa Zhang como motivación,
        sin rederivar su bloque de información pública.
 \end{itemize}
 \medskip
 Mi contribución será la corrección específica, sus condiciones y la
 comprobación reproducible. La revisión de literatura es delimitada y
 continuará antes del manuscrito final.
 \medskip
 \small
 Registro y alcance: \texttt{proposal/literature-search.md}.
\end{frame}

\begin{frame}[fragile]{Lean: el resultado algebraico de cuatro grupos}
 \small
 En notación económica: $\mathcal K(x)<\mathcal K(y)$.
 En Lean, \texttt{informationCoefficient} es exactamente la función
 racional definida en la propuesta.
 \begin{lstlisting}
theorem four_group_information_increases (a n3 n4 e3 e4 t R x y : (*@\(\mathbb R\)@*))
    (ha : 0 < a) (hn3 : 0 < n3) (hn4 : 0 < n4)
    (he3 : 0 < e3) (he4 : 0 < e4) (ht : 0 (*@\(\leq\)@*) t) (hR : 0 < R)
    (hxy : x < y) (hy1 : y (*@\(\leq\)@*) 1) :
    informationCoefficient a n3 n4 e3 e4 t R x <
      informationCoefficient a n3 n4 e3 e4 t R y := by
 \end{lstlisting}
 \small
 $n_j=\lambda_j/g_j$, $e_3=\sigma_\epsilon^2$,
 $e_4=\sigma_\epsilon^2+\tau_2^2$, $t=s$.
 \medskip
 \begin{itemize}
  \item Verifica la desigualdad a partir de hipótesis primitivas positivas.
  \item El contraejemplo también está probado, sin \texttt{sorry}.
  \item Falta formalizar la conexión con las esperanzas gaussianas.
        La prueba algebraica no certifica todo el modelo estocástico.
 \end{itemize}
\end{frame}

\begin{frame}{5. Plan y riesgos}
 \begin{enumerate}
  \item Derivar información condicional, las cuatro FOC y el equilibrio.
  \item Simular cambios de transparencia, riesgo y relevancia de grupos.
  \item Formalizar los resultados del manuscrito final y mapear sus hipótesis.
 \end{enumerate}
 \medskip
 \textbf{Riesgo principal:} si divulgar información también cambia el
 sentimiento, la comparación deja de mantener los pesos emocionales fijos.
 El signo exige una nueva derivación.
 \medskip
 \par
 Delimitaré la inferencia adicional desde precios y distinguiré
 sensibilidad, error de valoración y bienestar.
 \medskip
 \par
 La entrega final incluirá el apéndice manuscrito propio y la
 formalización de AppliedModelingLib del commit fijado del paper.
\end{frame}

\begin{frame}{La pregunta y la contribución que se evaluará}
 \large
 ¿Cuándo la transparencia fortalece la incorporación de
 fundamentos al precio?
 \vspace{0.4cm}
 \normalsize
 \begin{itemize}
  \item Corregiré la incertidumbre residual de la información pública.
  \item Probaré una respuesta esperada creciente con condiciones explícitas.
  \item Examinaré qué condición falla cuando cambia el sentimiento.
 \end{itemize}
 \medskip
 Incluso con información completa pueden persistir sesgos de expectativas.
 El resultado propuesto identifica una mejora en la sensibilidad informativa.
\end{frame}
\end{document}
